\section{注意力的 CUDA 实现}
\label{l12:sec:cuda}

\subsection{矩阵乘法流水线与存储布局}

注意力中的投影、分数计算和值聚合都能归结为通用矩阵乘法（General Matrix Multiplication，GEMM）。可以调用 CUDA 基础线性代数库（cuBLAS），也可以用分块矩阵乘法核函数实现。逐行的指数归一化穿插在两次注意力矩阵乘法之间，需要额外处理归约与因果性。

\begin{table}[htbp]
\centering
\caption{单序列注意力的基线计算步骤。多头分数和输出按各头分别计算。}
\label{l12:tab:pipeline}
\begin{tabularx}{\linewidth}{@{}L{2.4cm}L{6.5cm}Y@{}}
\toprule
步骤 & 运算与形状 & GPU 上的主要操作\\
\midrule
三个投影 & $(N\times D)(D\times3D)\to N\times3D$ & 可融合的矩阵乘法\\
分数计算 & $(N\times d)(d\times N)\to N\times N$ & 矩阵乘法与 $1/\sqrt d$ 缩放\\
因果归一化 & $S\to P$，保持 $N\times N$ & 行内最大值与求和归约\\
值聚合 & $(N\times N)(N\times d)\to N\times d$ & 矩阵乘法\\
头合并与投影 & $(N\times D)(D\times D)\to N\times D$ & 重排与矩阵乘法\\
\bottomrule
\end{tabularx}
\end{table}

缩放可以作为分数矩阵乘法的标量系数完成，从而避免再单独遍历 $N^2$ 个元素。数学上的 $K^{\mathsf T}$ 也不必先物化成一个新矩阵；支持转置操作的矩阵乘法接口可以按照原缓冲区的布局读取。但行主序与列主序、转置标志以及前导维度必须相互一致，才能实现式~\eqref{l12:eq:attention} 指定的乘法。

加入批量维后，输入常表示为 $[B,N,D]$。把融合投影的最后一维拆成 $[3,h,d]$，得到 $[B,N,3,h,d]$。抽取查询部分后得到 $[B,N,h,d]$，再交换头轴与序列轴，形成 $[B,h,N,d]$。分数矩阵则是 $[B,h,N,N]$。维度重塑（reshape）可以只改变同一线性存储的解释；轴置换是否需要数据搬运，取决于后续核函数是否能直接处理相应的步长。布局选择会影响合并访存和矩阵乘法接口的使用方式。

以下核函数只处理一个头的一张分数矩阵。假定它按行主序连续存储，因此 $S_{ij}$ 的地址偏移是 $iN+j$。输入 $S$ 已经包含 $1/\sqrt d$ 缩放，但尚未归一化。输出 $P$ 使用另一块同形状缓冲区，并将掩码位置明确置零，供后续稠密矩阵乘法使用。

\subsection{数值稳定的因果归一化}

直接计算 $\exp(S_{ij})$ 可能溢出。对同一行减去一个常数 $c_i$ 不改变归一化结果，因为分子和分母中的公共因子 $\exp(c_i)$ 可以约去：
\begin{equation}
 \frac{\exp(S_{ij})}{\sum_{k=0}^{i}\exp(S_{ik})}
 =\frac{\exp(S_{ij}-c_i)}
 {\sum_{k=0}^{i}\exp(S_{ik}-c_i)}.
 \label{l12:eq:shift-invariance}
\end{equation}
取有效列中的最大值 $m_i=\max_{0\leq j\leq i}S_{ij}$，所有参与计算的指数输入都不大于零，至少有一项指数等于一。定义行分母 $s_i$，得到
\begin{equation}
\begin{aligned}
 m_i&=\max_{0\leq j\leq i}S_{ij},\qquad
 s_i=\sum_{j=0}^{i}\exp(S_{ij}-m_i),\\
 P_{ij}&=\begin{cases}
 \exp(S_{ij}-m_i)/s_i,&0\leq j\leq i,\\
 0,&i<j<N.
 \end{cases}
\end{aligned}
\label{l12:eq:stable-causal}
\end{equation}
在有效分数有限的条件下，$1\leq s_i\leq i+1$。这既避免指数上溢，也使分母具有明确的正下界。非常小的权重仍可能发生浮点下溢，但它们对加权和的影响相应较小。

最大值与分母的范围必须遵守同一掩码。例如，在行 $i=1$，假设三个分数依次为 $1000$、$1001$ 和 $9000$。只有前两列可见，正确计算取 $m_1=1001$，使用 $(e^{-1},1)$ 归一化，得到约 $(0.268941,0.731059,0)$。若错误地把未来列的 $9000$ 纳入最大值，再只累加前两列，两项指数都可能下溢成零。若把未来列也纳入分母，则直接改变了因果分布。实现时遍历有效列，可以同时避免这两类问题。

这一步的并行结构不同于逐元素激活：同一行的每个输出都依赖该行共同的最大值和分母。因此必须先在行内归约，再把归约结果提供给所有负责该行的线程。

\subsection{一行一个线程块的计算组织}

令线程块包含 $B_{\mathrm t}$ 个线程，块编号 $i=\code{blockIdx.x}$ 对应行号，线程编号 $u=\code{threadIdx.x}$。线程 $u$ 按步长 $B_{\mathrm t}$ 访问 $j=u,u+B_{\mathrm t},u+2B_{\mathrm t},\ldots$。前两次遍历只处理 $j\leq i$；最后一次遍历覆盖 $j<N$，以写出完整的一行。

计算先由每个线程求自己负责部分的最大值，再在线程块内归约得到 $m_i$；随后，每个线程累加自己的 $\exp(S_{ij}-m_i)$，再次归约得到 $s_i$；最后，各线程独立写出概率。没有有效列的线程在最大值归约中贡献 $-\infty$，在求和归约中贡献零，仍然参与全部同步。这样，即使第一行只有一个有效元素，整个线程块的控制流程也保持一致。

例如取 $B_{\mathrm t}=256$，处理行 $i=700$ 时，可见列为 $0,\ldots,700$。线程 $0$ 负责列 $0,256,512$；线程 $189$ 负责 $189,445$，下一项 $701$ 已越过可见边界。因此，一行的长度可以大于线程块大小；增加长度表现为每个线程循环更多次。线程块数对应行数，也需要满足设备的网格维度限制。若把多个批次和多个头的行展平，则总行号 $r$ 对应的因果位置应取 $i=r\bmod N$，不能直接把 $r$ 当作掩码边界。

下面采用共享内存上的二叉树归约。设 $B_{\mathrm t}$ 为二的幂，线程先把局部结果写入共享数组；归约步长依次取 $B_{\mathrm t}/2,B_{\mathrm t}/4,\ldots,1$。每一步中，前半部分线程组合两个位置的结果，再由块内屏障保证下一步读到已完成的写入。归约结束后，共享数组的第零项包含整行需要的统计量。

\noindent\begin{minipage}{\linewidth}
\lstinputlisting[firstline=1,lastline=28,firstnumber=1,caption={线程块内的最大值或求和归约},label={l12:lst:reduce}]{code/causal_softmax.cu}
\end{minipage}

模板参数 \code{ReduceMax} 在编译时确定归约操作。每个线程都调用这个函数，包括局部结果为单位元的线程。循环结束时的屏障保证归约已经完成；随后先把 \code{scratch[0]} 读入每个线程自己的寄存器变量，再做一次屏障。这最后一次同步保证所有线程都已读取旧结果，后面的求和归约才可以安全地复用同一块共享内存。

核函数采用三次按行遍历。第二次遍历将指数临时写入输出缓冲区 $P$，第三次遍历再除以公共分母；这样每个有效位置只需计算一次指数。最后一次遍历使用相同的线程分工，因此每个线程归一化自己刚才写入的位置。$P$ 的上三角部分尚未初始化，条件表达式会直接写零，不读取其中的旧内容。

\noindent\begin{minipage}{\linewidth}
\lstinputlisting[firstline=30,lastline=63,firstnumber=30,caption={单头、单序列的因果归一化核函数},label={l12:lst:kernel}]{code/causal_softmax.cu}
\end{minipage}

这里使用 \code{std::size_t} 计算行偏移，使乘积 $iN$ 在较宽的无符号索引类型中形成。输入与输出缓冲区需要各自至少容纳 $N^2$ 个单精度数，且二者的存储区间互不重叠。有效输入分数应为有限值；仅用于表示未来位置的上三角值不会被读取。

暂存指数和重新计算指数是两种保持同一数学结果的实现方式。暂存版本减少指数函数调用，同时增加一次中间结果写入；重新计算版本可以在最后一轮再次读取 $S$ 并求指数。哪一种更快，取决于指数运算吞吐、显存带宽和缓存行为，需要按目标设备测量。上面的版本主要用于清楚展示归约与同步关系。

\subsection{启动条件与实现检查}

核函数通过一维网格启动 $N$ 个线程块，每块固定为 $256$ 个线程。启动之前应验证设备允许相应的网格大小、线程块大小和共享内存需求。共享数组只需 $256\times4=1024$ 字节，大小与 $N$ 无关。较长上下文增加线程内循环次数和总行数，并不要求共享内存容纳整行。

\noindent\begin{minipage}{\linewidth}
\lstinputlisting[firstline=65,lastline=90,firstnumber=65,caption={主机端的参数检查与异步启动},label={l12:lst:launch}]{code/causal_softmax.cu}
\end{minipage}

这个辅助函数返回参数、设备查询或启动阶段发现的错误码，不等待核函数执行结束。同一流中的后续矩阵乘法按提交顺序读取 $P$；需要确认执行完成或测量耗时时，还应在合适的位置等待该流或记录设备事件。函数无法从裸指针确定缓冲区的实际容量，调用者仍需满足前述容量与不重叠条件。

正确性检查可以直接使用第~\ref{l12:subsec:example} 节的分数矩阵。预期输出为式~\eqref{l12:eq:example-output} 中的 $P$；随后计算 $PV$ 应得到对应的 $O$。还应覆盖 $N=1$、长度小于线程数、长度不等于线程数的整数倍以及长度大于线程数的情况。检查的对象包括有效行和接近一、上三角严格为零、输出有限，以及与高精度逐行参考计算的一致性。单精度归约改变加法顺序，通常会产生很小的舍入差异，应根据数据类型比较数值误差。

性能上，一个线程块处理一行使行内通信局限于共享内存和块内屏障。相邻线程在同一次循环中访问相邻列，便于合并访存；早期行只有少量有效元素，部分线程参与同步但没有实际数据工作。更大的线程块可以减少长行的循环次数，也会增加线程和寄存器资源占用。线程数的选择需要同时考虑行长、活跃线程比例和设备并发能力。
